%%%%%%%%%%%%%%%%
%%% exod 22 %%%
%%%%%%%%%%%%%%%%

\Chapter{22}

\Verse{1}
{
	\hE{אִם בַּמַּחְתֶּרֶת יִמָּצֵא הַגַּנָּב וְהֻכָּה וָמֵת אֵין לוֹ דָּמִים׃}
}
{
	\eE{
		“If the thief will be found while breaking in and will
		be struck, and he dies, there is no blood for him.
	}
}
{
%	\footnote{
%		no blood for him. Meaning: there is no bloodguilt for killing him. The
%		person who killed the thief who was in the act of breaking in (at night)
%		cannot be held legally responsible for the death.
%	}
}

\Verse{2}
{
	\hE{אִם זָרְחָה הַשֶּׁמֶשׁ עָלָיו דָּמִים לוֹ שַׁלֵּם יְשַׁלֵּם אִם אֵין לוֹ וְנִמְכַּר בִּגְנֵבָתוֹ׃}
}
{
	\eE{
		If the sun has risen on him there is blood for him. He shall pay. If
		he does not have it, then he shall be sold for his theft.
	}
}
{
%	\footnote{
%		22.2. If the sun has risen on him there is blood for him. If the
%		theft was in the daytime—which is less frightening and when escape,
%		help, and witnesses are more possible—then one is legally
%		responsible if one kills the thief.
%
%		Footnote 22:2. He shall pay. If the thief is caught, not killed, he
%		must pay for his theft.
%
%		Footnote 22:2. If he does not have it, then he shall be sold for his
%		theft. If the thief does not have the money to pay for his theft,
%		then he is sold into servitude.
%	}
}

\Verse{3}
{
	\hE{אִם הִמָּצֵא תִמָּצֵא בְיָדוֹ הַגְּנֵבָה מִשּׁוֹר עַד חֲמוֹר עַד שֶׂה חַיִּים שְׁנַיִם יְשַׁלֵּם׃}
}
{
	\eE{
		If the theft will be found in his hand—from ox to ass to
		sheep—alive, he shall pay two.
	}
}
{
}

\Verse{4}
{
	\hE{כִּי יַבְעֶר אִישׁ שָׂדֶה אוֹ כֶרֶם וְשִׁלַּח אֶת בְּעִירֹה וּבִעֵר בִּשְׂדֵה אַחֵר מֵיטַב שָׂדֵהוּ וּמֵיטַב כַּרְמוֹ יְשַׁלֵּם׃}
}
{
	\eE{
		“If a man will have a field or a vineyard grazed, and he
		will let his beast go and graze in another person’s field,
		he shall pay the best of his field and the best of his
		vineyard.
	}
}
{
}

\Verse{5}
{
	\hE{כִּי תֵצֵא אֵשׁ וּמָצְאָה קֹצִים וְנֶאֱכַל גָּדִישׁ אוֹ הַקָּמָה אוֹ הַשָּׂדֶה שַׁלֵּם יְשַׁלֵּם הַמַּבְעִר אֶת הַבְּעֵרָה׃}
}
{
	\eE{
		“If fire will break out and find thorns and stacked grain
		or standing grain or a field is consumed, the one who set
		the blaze shall pay.
	}
}
{
}

\Verse{6}
{
	\hE{כִּי יִתֵּן אִישׁ אֶל רֵעֵהוּ כֶּסֶף אוֹ כֵלִים לִשְׁמֹר וְגֻנַּב מִבֵּית הָאִישׁ אִם יִמָּצֵא הַגַּנָּב יְשַׁלֵּם שְׁנָיִם׃}
}
{
	\eE{
		“If a man will give his neighbor money or items to watch,
		and it will be stolen from the man’s house, if the thief
		will be found he shall pay two.
	}
}
{
}

\Verse{7}
{
	\hE{אִם לֹא יִמָּצֵא הַגַּנָּב וְנִקְרַב בַּעַל הַבַּיִת אֶל הָאֱלֹהִים אִם לֹא שָׁלַח יָדוֹ בִּמְלֶאכֶת רֵעֵהוּ׃}
}
{
	\eE{
		If the thief will not be found, then the owner of the
		house shall be brought near to {\God}: that he has not put
		out his hand into his neighbor’s property.
	}
}
{
%	\footnote{
%		brought near to God: that he has not put out his hand. The wording
%		(starting with ’im) plus the notice of being brought “near to God”
%		indicate that he is to take an oath that he has not taken the missing
%		property for which he was responsible. The wording occurs again in 22:10
%		with an explicit reference to an oath.
%	}
}

\Verse{8}
{
	\hE{עַל כּל דְּבַר פֶּשַׁע עַל שׁוֹר עַל חֲמוֹר עַל שֶׂה עַל שַׂלְמָה עַל כּל אֲבֵדָה אֲשֶׁר יֹאמַר כִּי הוּא זֶה עַד הָאֱלֹהִים יָבֹא דְּבַר שְׁנֵיהֶם אֲשֶׁר יַרְשִׁיעֻן אֱלֹהִים יְשַׁלֵּם שְׁנַיִם לְרֵעֵהוּ׃}
}
{
	\eE{
		“Over every case of an offense, over an ox, over an ass,
		over a sheep, over clothing, over every loss about which
		one will say that ‘This is it,’ the word of the two of
		them shall come to {\God}. The one whom {\God} will implicate
		shall pay two to his neighbor.
	}
}
{
%	\footnote{
%		implicate. The verb here is plural. Some therefore take its subject,
%		’lhîm, not to refer to God here and in the preceding verse. (See Rashi
%		and Ibn Ezra.) Rather, it is understood to refer to the judges before
%		whom the case is brought.
%	}
}

\Verse{9}
{
	\hE{כִּי יִתֵּן אִישׁ אֶל רֵעֵהוּ חֲמוֹר אוֹ שׁוֹר אוֹ שֶׂה וְכל בְּהֵמָה לִשְׁמֹר וּמֵת אוֹ נִשְׁבַּר אוֹ נִשְׁבָּה אֵין רֹאֶה׃}
}
{
	\eE{
		“If a man will give an ass or an ox or a sheep or any
		animal to his neighbor to watch, and it dies or is injured
		or is seized—no one seeing—
	}
}
{
}

\Verse{10}
{
	\hE{שְׁבֻעַת יְהֹוָה תִּהְיֶה בֵּין שְׁנֵיהֶם אִם לֹא שָׁלַח יָדוֹ בִּמְלֶאכֶת רֵעֵהוּ וְלָקַח בְּעָלָיו וְלֹא יְשַׁלֵּם׃}
}
{
	\eE{
		an oath of YHWH shall be between the two of them: that he
		has not put out his hand into his neighbor’s property; and
		its owner shall take it, and he shall not pay.
	}
}
{
%	\footnote{
%		he shall not pay. The person who was watching the animal does not pay
%		anything to the animal’s owner.
%	}
}

\Verse{11}
{
	\hE{וְאִם גָּנֹב יִגָּנֵב מֵעִמּוֹ יְשַׁלֵּם לִבְעָלָיו׃}
}
{
	\eE{And if it will be stolen from him, he shall pay its owner.}
}
{
}

\Verse{12}
{
	\hE{אִם טָרֹף יִטָּרֵף יְבִאֵהוּ עֵד הַטְּרֵפָה לֹא יְשַׁלֵּם׃}
}
{
	\eE{
		If it will be torn, he shall bring it in witness; he shall
		not pay for the torn one.
	}
}
{
%	\footnote{
%		witness. Rashi understands this to mean that he brings witnesses. Ibn
%		Ezra suggests that it means that he brings a portion of the torn animal
%		as evidence. The latter is the understanding of most contemporary
%		translators.
%	}
}

\Verse{13}
{
	\hE{וְכִי יִשְׁאַל אִישׁ מֵעִם רֵעֵהוּ וְנִשְׁבַּר אוֹ מֵת בְּעָלָיו אֵין עִמּוֹ שַׁלֵּם יְשַׁלֵּם׃}
}
{
	\eE{
		And if a man will ask it of his neighbor, and it is
		injured or dies, its owner not being with it, he shall
		pay.
	}
}
{
%	\footnote{
%		ask. That is, borrow.
%	}
}

\Verse{14}
{
	\hE{אִם בְּעָלָיו עִמּוֹ לֹא יְשַׁלֵּם אִם שָׂכִיר הוּא בָּא בִּשְׂכָרוֹ׃}
}
{
	\eE{
		If its owner is with it, he shall not pay. If it was
		hired, he has its hire coming to him.
	}
}
{
}

\Verse{15}
{
	\hE{וְכִי יְפַתֶּה אִישׁ בְּתוּלָה אֲשֶׁר לֹא אֹרָשָׂה וְשָׁכַב עִמָּהּ מָהֹר יִמְהָרֶנָּה לּוֹ לְאִשָּׁה׃}
}
{
	\eE{
		“And if a man will deceive a virgin who is not betrothed,
		and he lies with her, he shall espouse her by a
		bride-price to him as a wife.
	}
}
{
}

\Verse{16}
{
	\hE{אִם מָאֵן יְמָאֵן אָבִיהָ לְתִתָּהּ לוֹ כֶּסֶף יִשְׁקֹל כְּמֹהַר הַבְּתוּלֹת׃}
}
{
	\eE{
		If her father will refuse to give her to him, he shall
		weigh out money corresponding to the bride-price of
		virgins.
	}
}
{
}

\Verse{17}
{
	\hE{מְכַשֵּׁפָה לֹא תְחַיֶּה׃}
}
{
	\eE{“You shall not let a witch live.}
}
{
%	\footnote{
%		a witch. Why does the text condemn only women and not male sorcerers?
%		Males are identified as practicing magic as well, as in the case of the
%		Egyptian magicians (Exod 7:11). The feminine form of this word occurs
%		only here out of the entire Tanak. And the Septuagint has the masculine
%		plural (which is used to include both men and women) here instead of the
%		Masoretic Text’s feminine singular. This suggests that the original text
%		forbids the practice of magic by either men or women (sorcerers or
%		witches). Note also: the text does not deny that magic works. (On the
%		contrary, the Egyptian magicians can turn staffs into snakes.) It
%		forbids the practice of magic.
%	}
}

\Verse{18}
{
	\hE{כּל שֹׁכֵב עִם בְּהֵמָה מוֹת יוּמָת׃}
}
{
	\eE{“Anyone who lies with an animal shall be put to death.}
}
{
}

\Verse{19}
{
	\hE{זֹבֵחַ לָאֱלֹהִים יחֳרָם בִּלְתִּי לַיהֹוָה לְבַדּוֹ׃}
}
{
	\eE{
		“One who sacrifices to gods shall be completely
		destroyed—except to YHWH alone.
	}
}
{
%	\footnote{
%		completely destroyed. Hebrew erem refers to the rule of destroying an
%		individual or community along with all his family, animals, and
%		possessions and thus dedicating them to the deity rather than profiting
%		in any way from them.
%	}
}

\Verse{20}
{
	\hE{וְגֵר לֹא תוֹנֶה וְלֹא תִלְחָצֶנּוּ כִּי גֵרִים הֱיִיתֶם בְּאֶרֶץ מִצְרָיִם׃}
}
{
	\eE{
		“And you shall not persecute an alien, and you shall not
		oppress him, because you were aliens in the land of Egypt.
	}
}
{
}

\Verse{21}
{
	\hE{כּל אַלְמָנָה וְיָתוֹם לֹא תְעַנּוּן׃}
}
{
	\eE{“You shall not degrade any widow or orphan.}
}
{
}

\Verse{22}
{
	\hE{אִם עַנֵּה תְעַנֶּה אֹתוֹ כִּי אִם צָעֹק יִצְעַק אֵלַי שָׁמֹעַ אֶשְׁמַע צַעֲקָתוֹ׃}
}
{
	\eE{
		If you will degrade them, when they will cry out to me
		I’ll hear their cry,
	}
}
{
}

\Verse{23}
{
	\hE{וְחָרָה אַפִּי וְהָרַגְתִּי אֶתְכֶם בֶּחָרֶב וְהָיוּ נְשֵׁיכֶם אַלְמָנוֹת וּבְנֵיכֶם יְתֹמִים׃}
}
{
	\eE{
		and my anger will flare, and I’ll kill you with a sword,
		and your wives will be widows and your children orphans!
	}
}
{
}

\Verse{24}
{
	\hE{אִם כֶּסֶף תַּלְוֶה אֶת עַמִּי אֶת הֶעָנִי עִמָּךְ לֹא תִהְיֶה לוֹ כְּנֹשֶׁה לֹא תְשִׂימוּן עָלָיו נֶשֶׁךְ׃}
}
{
	\eE{
		“If you will lend money to my people, to the poor who is
		with you, you shall not be like a creditor to him; you
		shall not impose interest on him.
	}
}
{
}

\Verse{25}
{
	\hE{אִם חָבֹל תַּחְבֹּל שַׂלְמַת רֵעֶךָ עַד בֹּא הַשֶּׁמֶשׁ תְּשִׁיבֶנּוּ לוֹ׃}
}
{
	\eE{
		If you take your neighbor’s clothing as security, you
		shall give it back to him by the sunset,
	}
}
{
}

\Verse{26}
{
	\hE{כִּי הִוא כְסוּתֹה לְבַדָּהּ הִוא שִׂמְלָתוֹ לְעֹרוֹ בַּמֶּה יִשְׁכָּב וְהָיָה כִּי יִצְעַק אֵלַי וְשָׁמַעְתִּי כִּי חַנּוּן אָנִי׃}
}
{
	\eE{
		because it is his only apparel; it is his clothing for his
		skin. In what will he sleep? And it will be, when he will
		cry out to me, that I shall listen, because I am gracious.
	}
}
{
}

\Verse{27}
{
	\hE{אֱלֹהִים לֹא תְקַלֵּל וְנָשִׂיא בְעַמְּךָ לֹא תָאֹר׃}
}
{
	\eE{
		“You shall not blaspheme {\God}, and you shall not curse a
		chieftain among your people.
	}
}
{
%	\footnote{
%		chieftain. Hebrew nî’. Elsewhere in the Torah the text refers to the
%		chieftain of a tribe (Num 2:3–29) or of a priestly family (Num 3:24–35)
%		or of a city (Gen 34:2). Here it may allude to a king of Israel, “a
%		chieftain among your people” (as in 1 Kings 11:34; Ezek 34:24), or it
%		may mean any of these leaders.
%	}
}

\Verse{28}
{
	\hE{מְלֵאָתְךָ וְדִמְעֲךָ לֹא תְאַחֵר בְּכוֹר בָּנֶיךָ תִּתֶּן לִי׃}
}
{
	\eE{
		“You shall not delay your fulfillment and your flowing.
		“You shall give me the firstborn of your sons.
	}
}
{
%	\footnote{
%		fulfillment and your flowing. Everyone is unsure of what this means.
%	}
}

\Verse{29}
{
	\hE{כֵּן תַּעֲשֶׂה לְשֹׁרְךָ לְצֹאנֶךָ שִׁבְעַת יָמִים יִהְיֶה עִם אִמּוֹ בַּיּוֹם הַשְּׁמִינִי תִּתְּנוֹ לִי׃}
}
{
	\eE{
		“You shall do this to your ox and to your sheep: Seven
		days it will be with its mother. On the eighth day you
		shall give it to me.
	}
}
{
}

\Verse{30}
{
	\hE{וְאַנְשֵׁי קֹדֶשׁ תִּהְיוּן לִי וּבָשָׂר בַּשָּׂדֶה טְרֵפָה לֹא תֹאכֵלוּ לַכֶּלֶב תַּשְׁלִכוּן אֹתוֹ׃}
}
{
	\eE{
		“And you shall be holy people to me. “And you shall not
		eat meat in the field that is torn. You shall throw it to
		the dog.
	}
}
{
%	\footnote{
%		meat that is torn. This refers to the meat of an animal that was killed
%		by another animal rather than having been killed in the sacred manner
%		prescribed by the law. The word for “torn” (Hebrew trph) later came to
%		be used for any prohibited foods.
%	}
}
